\section{Introducción y contexto de los datos}

En un material magnético, cada átomo se comporta como una pequeña brújula, un espín $\mathbf S_i$ de norma uno que puede apuntar en cualquier dirección del espacio. A alta temperatura las brújulas fluctúan sin orden; al enfriar, sus interacciones las alinean y aparece una textura ordenada. Cuando el material carece de simetría de inversión, la interacción de Dzyaloshinskii y Moriya \citep{dzyaloshinskii1958,moriya1960}, de intensidad $\kdm$, favorece que los espines vecinos giren un ángulo fijo. El resultado son espirales y, en presencia de un campo externo $\Hex$, texturas de tipo skyrmión con carga topológica no necesariamente entera en un disco de bordes libres \citep{nagaosa2013}. El objeto de estudio es un disco de radio $R_d=18.5$ sitios y cinco capas, con $N=5425$ espines sobre una red cúbica de $39\times39\times5$ posiciones. El modelo y el esquema de simulación provienen del trabajo de J.~D.~Agudelo-Giraldo y colaboradores: un hamiltoniano de Heisenberg extendido para nanodiscos, con intercambio, interacción de Dzyaloshinskii--Moriya, anisotropía magnetocristalina y energía de Zeeman, muestreado con Metropolis--Hastings acelerado en GPU y enfriamiento progresivo \citep{mendezrondon2026}. Ese grupo ha aplicado la simulación Monte Carlo atomística a películas delgadas nanocristalinas, nanotubos y películas ultradelgadas \citep{agudelo2016,agudelo2017,agudelo2020,agudelo2023}.

Para un estadístico, el problema tiene una lectura directa. Una configuración $\mathbf s=(\mathbf S_1,\dots,\mathbf S_N)$ es un vector aleatorio de dimensión $2N=10\,850$ con densidad de Boltzmann; las cantidades físicas son esperanzas bajo esa densidad, y su constante de normalización es una integral imposible de calcular. Por eso se recurre a una cadena de Markov (Metropolis) cuya distribución estacionaria es la de Boltzmann. Las dos preguntas del proyecto son estadísticas: ¿a qué temperatura cambia de régimen el disco y con qué incertidumbre se conoce?, y ¿alcanzaron las cadenas su distribución estacionaria con el número de pasos simulado?

El diseño simula cada estado de parámetros con 100 copias independientes. Una copia arranca a 20~K, se enfría en pasos de 0.1~K hasta 6~K y en cada temperatura continúa desde la configuración que dejó la anterior. Cada copia es, entonces, un camino de enfriamiento distinto y todos son igualmente válidos. A temperatura alta los caminos deben describir la misma distribución; a temperatura baja pueden quedar en estados finales diferentes, porque el paisaje de energía tiene varios valles casi equivalentes. \citet{rincon2026} documentan este comportamiento en nanopartículas de cobre enfriadas a distintas tasas: las curvas de energía coinciden a temperatura alta, divergen levemente a temperatura baja y el pico del calor específico identifica las transiciones. Este informe cuantifica esa divergencia con herramientas del curso, distingue diversidad del protocolo y falta de mezcla de equilibrio, y evalúa si aumentar el número de barridos puede mejorar el muestreo. Las realizaciones son válidas para el protocolo; su representatividad de Boltzmann exige evidencia adicional.

\begin{table}[H]
\centering\small
\caption{Estados simulados. Intercambio $J=1$ y anisotropía magnetocristalina cúbica $\Kan=0.05$ en todos los casos; el campo externo $\Hex$ es perpendicular al disco ($\gamma=90^\circ$). Cada lote de Kaggle ejecutó dos estados en dos GPU T4.}
\label{tab:estados}
\begin{tabular}{lcccccccc}
\toprule
Estado & 1 & 2 & 3 & 4 & 5 & 6 & 7 & 8\\
\midrule
$\kdm$ & 0.05 & 0.15 & 0.33 & 0.60 & 1.00 & 0.30 & 0.60 & 0.75\\
$\Hex$ & 0 & 0 & 0 & 0 & 0 & 0.05 & 0.10 & 0.30\\
Lote & 1 & 1 & 2 & 2 & 3 & 3 & 4 & 4\\
\bottomrule
\end{tabular}
\end{table}

Los datos son series temporales simuladas, no observadas. Por cada estado, temperatura y copia se registraron, en cada uno de los 10\,000 barridos de medición, la energía total $E$ normalizada por $E_0$, sus cuatro componentes (intercambio, anisotropía, Dzyaloshinskii y Moriya, campo) y la norma de la magnetización por espín $m$. Las 141 temperaturas por estado y los ocho estados suman 1128 matrices de $10\,000\times100$ por variable, unos 25~GB en formato comprimido. No hay valores faltantes; el análisis no requiere el tratamiento de ceros que la guía exige a las series financieras, pues $E$ y $m$ son continuas y $m>0$ casi seguramente.

\subsection{Preguntas e hipótesis operativas}

El estudio pregunta dónde se sitúa el máximo energético, si su posición se asocia con $\kdm$ y $\Hex$, y en qué regiones los caminos mezclan con el presupuesto disponible. La hipótesis de asociación nula entre $\kdm$ y $T^*$ se evalúa con Spearman y permutación exacta. Su no rechazo no establece independencia. La compatibilidad entre caminos se examina con $\Rhat$, Geweke, MSER-5 y homogeneidad; pasar esos diagnósticos no demuestra exploración global.

La hipótesis física de trabajo relaciona la escala de las texturas con el costo del movimiento local. El factor de estructura, las cargas finales y las mezclas de magnetización evalúan aspectos complementarios. Para separar efectos de campo e interacción quiral haría falta un diseño factorial. Donde no se demuestra equilibrio, las frecuencias finales se interpretan como resultados del protocolo de enfriamiento.
